\subsection{Retrieval-Augmented Generation}

RAG is used add a retrieval step in an LLM workflow. The model no longer relies on training data, it first retrieves through an external collection all the documents it might need, then select chunks of it. It will select a passage or a record and put it into context before letting the model produce an answer based on the collected data \cite{lewis2020rag}. Recent surveys mention this as the new rising technique that allows LLMs to utilise the most out of the user data to give the best answer possible in retrieval-augmented settings \cite{gao2024ragsurvey,huang2024ragsurvey}. This explains our reasoning for tracing every step of the experiment: the data dictates what can be recommended, the index show how we can parse the data that retriever uses to decide which candidates to add to the list, Context decides how the list gets presented to the LLM which then generates an output that our evaluator uses to make a decision on the outcome. 

For this thesis, the important point is not that every recommender is a textbook RAG system. The useful connection is narrower. When a recommendation pipeline stores movie titles, genres, synopses, tags, user preferences, or reviews as searchable text, the retrieved item record can become the evidence used for ranking, reranking, or explanation. In that setting, recommendation still returns a ranked list rather than a prose answer, but the dependence on retrieved text is similar. The candidate set is shaped before the final ranking decision. If an LLM is later used to rerank items or explain them, the retrieved fields become model-facing context. A recommendation can then change because a different record was retrieved, because the same record was represented differently, or because the model interpreted the retrieved text differently. Literature on LLM-powered recommendation and LLM-based ranking already treats language models as components that can represent preferences, process natural-language item descriptions, and order candidate lists \cite{wu2024llmrecsurvey,zhao2024llmrecsystems,hou2024llmrankers,sun2024rankgpt}.

It is important to mention that not all recommender systems are RAG systems, but in a context like this a retrieval tool only makes sense. A recommendation pipeline makes a decision based on a users preferences, age, past reviews, but also based on metadata specific to each item in its dataset and uses a mix of all of this data to do its final ranking and explanation. When the LLM is later used to rerank the items or give an explanation based on them, the retrieval data becomes the primary source of truth in this kind of pipeline. The recommendation stops being static too, it can change base on the items that were retrieved, because the data was presented in a different way or because the models interpretation wasn't the same. Studies about LLM-based ranking believes LLMs are components that can take preferences in consideration, process items description even as a natural sounding text and order a candidate list based on it \cite{wu2024llmrecsurvey,zhao2024llmrecsystems,hou2024llmrankers,sun2024rankgpt}.

This are the assumptions RAGPoison takes in consideration, it treats retrieval, ranking, metric computation and the collection of traces as separate part of the same benchmark, instead of a simple model call \cite{ragpoisonlab_repo}. 

\subsection{LLM-Powered Recommender Systems}

Large language models enter recommender systems in several roles. Surveys describe their use for preference modeling, natural-language interaction, explanation generation, item representation, sequential recommendation, and ranking or reranking \cite{wu2024llmrecsurvey,zhao2024llmrecsystems}. In many practical designs, however, the LLM does not replace the entire recommender. A conventional retrieval or scoring stage first narrows the catalog. The model then receives a shorter candidate list, often serialized as item titles, genres, summaries, or other textual fields. This division is attractive because it keeps the expensive model call small and keeps a familiar retrieval pipeline in place. It also means that earlier retrieval choices decide what the LLM ever sees. If a candidate is not retrieved, no later reranker can promote it. If a poisoned candidate is retrieved, the LLM may be asked to reason over it as if it were ordinary item metadata.

LLMs can have many roles in a recommender system, Surveys actually mention their usefulness for natural language interactions, generating explanations, recommendation and reranking \cite{wu2024llmrecsurvey,zhao2024llmrecsystems}. But unfortunately, an LLM is still unable to replace and entire recommender system. It will still need the starter data to rank, When it receives a small pool list it then can rank them or conversate on them. this keeps the model call cheap and avoid it losing context due too much shared data, it also helps with have a stable pipeline in place. It means that the LLM will only get as much context as its shared with it from the previous retrieval stage. If a poisoned candidate is is retrieved alongside the others the LLM might be asked to change its position in the list against the users best interest.  

Reranking is the role most relevant to this thesis. Work on LLM recommenders and rankers shows that large models can be prompted to order candidate items or documents without being trained as a traditional recommender for that exact task \cite{hou2024llmrankers,sun2024rankgpt}. The usual input is not just a movie identifier. It is a small piece of text that describes each candidate and asks the model to produce an order. That changes the threat model. Candidate text that once looked like passive metadata can become prompt-visible input. A title, synopsis, tag list, or injected field may sit next to trusted instructions in the same model call. The model is then expected to treat those fields as data, not as instructions. Indirect prompt-injection research shows why that boundary is fragile when applications consume external text \cite{greshake2023indirectpi,yi2025bipia,wallace2024instructionhierarchy}.

Reranking is the LLM feature that is most relevant for our usecase. Previous work on LLM recommenders show that a model can be prompted to rank a list without ever being trained as a recommender \cite{hou2024llmrankers,sun2024rankgpt}. The input would contain a small piece of text from each candidate and then ask the model to order them. This makes the usually considered passive data that come with a movie relevant since it's included in the prompt. The model then treats the data into the expected output. Some studies  raise the issue of indirect prompt injection where external data inadvertently affects the models actions \cite{greshake2023indirectpi,yi2025bipia,wallace2024instructionhierarchy}.

This point should be separated from broad claims that discuss the unsafety of LLMs. It should be discussed in the scope of our thesis through two points : how the item is retrieved into the candidate set and how the LLM interprets it. LLM ranking is another layer where the recommendation can be influenced due to wording order or payload instructions. 

\subsection{Retrieval and Ranking Metrics}

Having multiple metrics for each aspect of ranking is essential for this study. Hit Rate (HR) at $k$ checks if at least one user relevant item is in the top-$k$ list. It is simple and easy to understand but doesnt give a result based on the position of the item only its presence is relevant. Mean reciprocal rank (MRR) adds what we were lacking in HR previously, giving value to an items position on the list. It gives more credit when to the first user relevant item appears in the early ranks. NDCG or Normalized discounted cumulative gain is similar to MRR, giving value to an times rank except it applies that rule to all user relevant items, and not just the first. Studies on recommenders discussed in detail how a single facing metrics that can capture all actions in a system doesn't exist, so each question should be backed by it's own metric \cite{herlocker2004evaluating}.

But after worrying about the quality of the recommender we should now ask a different question: did the attack reach its target? In this thesis ASR or attack success rate is the metric we use to answer that question when applicable. In target-oriented attacks if target item reaches the top-$k$ ASR gives it a value based on its rank. This metric is different from others as it doesnt worry about the quality of the baseline list, which is relevant especially for targeted promotion and prompt injection. It solves the issue where a system might seem stable on the other metrics and yet a target movie was pushed in the top-$k$ without the other metrics accounting for it.  

Using all these metrics together prevents us from ignoring any specific part of our system. A poisoned index could push a single item to top-$k$ without ruining the recommendation quality and keeping the other metrics value similar to the baseline, while another attack could completely  reduce the quality of all targets. reporting the quality and the success of the attack is thus primordial to give a complete picture of the experiments ahead

\subsection{Recommender-System Poisoning}

Poisoning attempts on recommender systems are older than the ones on current RAG systems, Early studies discussed how modified rating data can influence the recommender \cite{li2016factorizationpoison} while later attacks targeted the top-$N$ outputs and used influence estimate to move items in a ranked list \cite{fang2020influencepoison}. Surveys analyzed and organized these studies based on the attacker goals, their injection points and defense strategy \cite{nguyen2024manipulatingrecsys,wang2024recpoisonsurvey}.

These studies are relevant for RAGPoison even though the injection point is different. A traditional poisoning behavior would change ratings or user profiles while we study controlled changes in document or their metadata before and after retrieval or ranking occurs. Nonetheless, the goal is still the same, promote a target, degrade general quality or modify the exposure of an item in top-$k$ lists. What changes is the way we reach the goal, A poisoned synopsis wouldn't look like modified data unless accounted for, an injected payload wouldnt change how a document is filtered in the matrix and yet, they will matter if the recommender retrieves those texts and passes them up to model for ranking. Recommender poisoning studies gives us that the vocabulary for targets, exposure and their countermeasures, while RAG security literature the document behavior \cite{jia2023pore,nguyen2024manipulatingrecsys,zhong2023retrievalcorpora,zou2025poisonedrag_usenix}.

These differences in attacks affects the defense mechanism since a detector specifically made to find fake users could miss poisoned metadata, training a model against attacks would be useful if it wasn't for the ability to tamper with retrieval indices and defense against prompt injection might reduce the risk of following payload instructions but would still fail if the item already got sent to it as a candidate. RAGPoison will use the older poisoning techniques as measurement while placing them in a retrieval setting, making the benchmark a bridge between two methods rather than replacing any of them \cite{ragpoisonlab_repo}. We will explore established concept and making a benchmark for them to later thing of defense mechanism against them.

\subsection{RAG Poisoning and Prompt Injection}

RAG poisoning studies attacks that compromise the material retrieved by a RAG system. One line of work inserts adversarial passages into retrieval corpora so that the system later retrieves and uses them \cite{zhong2023retrievalcorpora}. PoisonedRAG frames the problem as knowledge corruption: faulty retrieved documents can shape generated outputs even without access to model parameters \cite{zou2025poisonedrag_usenix}. Other work studies jailbreak-oriented poisoning, stealthy or human-imperceptible changes, backdoor-like triggers, and practical poisoning settings \cite{deng2024pandora,zhang2024humanimperceptible,xue2024badrag,zhang2026corruptrag}. Benchmark and detection work further shows that RAG poisoning should be evaluated as a system behavior, not only as a single bad answer \cite{zhang2025ragpoisonbench,tan2025revprag}. The shared lesson is that retrieved data is not neutral once downstream components act on it.

Prompt injection is related but not identical. Direct prompt injection comes from user-supplied prompt text. Indirect prompt injection enters through external material that the application retrieves, reads, or observes on the user's behalf \cite{greshake2023indirectpi}. Benchmarks and defenses study how such outside inputs can conflict with higher-priority instructions, and they propose strategies such as structured separation, spotlighting untrusted content, and instruction hierarchies \cite{yi2025bipia,debenedetti2024agentdojo,chen2024struq,hines2024spotlighting,wallace2024instructionhierarchy}. Retrieval poisoning asks how manipulated records enter the candidate set or context. Prompt injection asks whether text inside that context can be interpreted as an instruction. A recommender affected by both can experience ranking movement, target exposure, or unsafe model behavior through the same retrieved item record.

The overlap is the main security motivation for this thesis. RAG poisoning explains how malicious or misleading content can reach the model-facing context. Indirect prompt injection explains why content that looks like data may still interfere with LLM behavior. Recommender poisoning explains why small changes in exposure, top-$k$ membership, or ranking position matter. Recent work has started to examine poisoning attacks on RAG in recommender systems directly \cite{nazary2025poisonrag}. RAGPoison Lab is positioned in this intersection as a bounded local benchmark: it separates clean and poisoned recommendation paths, keeps attack descriptions sanitized, and records evidence needed to compare behavior without claiming to represent every production recommender \cite{ragpoisonlab_repo}. The thesis therefore treats attack-family labels as experimental categories. Retrieval poisoning concerns how records enter the candidate set or context. Prompt-injection-style poisoning concerns how text in that context may challenge instruction boundaries. Recommender poisoning concerns exposure and ranking effects. Keeping those meanings separate makes the later results easier to interpret.

\subsection{Red-Teaming and Benchmark Design}

Red-teaming is an attempt to find the unsafe and unwanted behavior before getting to deployment. And in the case of LLMs this usually involves test generation, automated adversarial evaluation, and benchmarks \cite{perez2022redteaming,mazeika2024harmbench}. Some studies share that scenario based, multi metric and with transparent evaluation benchmarks avoid the mistake of isolated examples in which it is easy to over interpret the results \cite{liang2023helm}. This means that in the context of this thesis a meaningful results will require the scenario, attack type, system configuration and other artifacts to be deeply documented. That will be the only way to truly understand and support our future claims.

We will also need to cater to the Rag and recommender systems own requirement in our project since a RAG benchmark not only needs to specify the language model but also its corpus and retrieval process. A recommender benchmark on the other hand needs to specify the users, items and metrics and a poisoning benchmark needs to specify the attackers goal, the scope of poisoning and a clear separation between the clean and attacked conditions. Without clear logging of all of these requirement, a change the recommendation could be coming from any of these steps without clearly knowing which one it was \cite{thakur2021beir,liang2023helm,zhang2025ragpoisonbench}. This is what will be achieved with RAGPoison, a smaller benchmark compared to live production systems but still useful since all steps between baseline and poisoned are auditable.



\subsection{Positioning Against Similar Benchmarks}

The related work also helps position this thesis against nearby benchmarks without turning the comparison into a leaderboard. Poison-RAG is the closest external reference because it studies adversarial data poisoning in retrieval-augmented recommender systems and also focuses on metadata such as tags and item descriptions \cite{nazary2025poisonrag}. PoisonedRAG and the broader RAG poisoning benchmark literature are close in the retrieval-security sense, but their main evaluation target is usually question answering or generation behavior rather than top-$k$ recommendation exposure \cite{zou2025poisonedrag_usenix,zhang2025ragpoisonbench}. Indirect prompt-injection benchmarks are also relevant because they study external text that reaches an LLM through an application workflow, which is the same class of risk created when a recommender passes item text to an LLM reranker \cite{yi2025bipia}. LLM recommendation benchmarks and zero-shot LLM ranker studies justify the ranking setting, but they are not poisoning benchmarks by themselves \cite{hou2024llmrankers,liu2023llmrecbench,geng2023p5}.

\begin{table}[H]
\centering
\small
\caption{Benchmark-positioning comparison for this thesis.}
\label{tab:benchmark-positioning-comparison}
\begin{tabular}{p{0.19\textwidth}p{0.18\textwidth}p{0.20\textwidth}p{0.33\textwidth}}
\toprule
Work & Main domain & Main attack or evaluation surface & Relation to RAGPoison Lab \\
\midrule
Poison-RAG & RAG recommender & Item metadata poisoning in a recommender setting & Closest comparison because it directly studies RAG-based recommender poisoning with item metadata. \\
PoisonedRAG & RAG question answering & Malicious knowledge injected into retrieval corpora & Close for retrieval poisoning, but the task is generation correctness instead of recommendation exposure. \\
RAG poisoning benchmark & RAG benchmark design & Multiple poisoning attacks and defenses across QA datasets & Useful for benchmark structure and threat-model comparison, not for direct metric comparison. \\
Indirect prompt-injection benchmarks & LLM-integrated applications & External content that conflicts with application instructions & Useful for the prompt-injection family and for explaining why retrieved text should be treated as untrusted. \\
LLM recommender and ranker studies & Recommendation and ranking & LLMs used as recommenders or rerankers & Useful for justifying the LLM reranking setting but not a poisoning baseline. \\
Traditional recommender poisoning & Collaborative filtering and top-$N$ recommendation & Fake ratings, fake users, or graph/profile manipulation & Useful for attacker goals and exposure concepts, while the injection point differs from metadata and retrieval poisoning. \\
\bottomrule
\end{tabular}
\end{table}

This thesis should therefore compare itself to earlier work at the level of setting, threat model, artifact design, and metric choice. A direct numerical comparison would be misleading unless the compared papers used the same dataset split, retrieval method, ranking method, poisoning budget, target definition, and metric implementation. The safer comparison is that RAGPoison Lab combines three ideas that are often studied separately: recommender poisoning goals, RAG-style retrieval poisoning surfaces, and prompt-visible LLM reranking. This makes the thesis narrower than a production security benchmark but more specific than a general LLM red-team prompt set.
